\resizebox{1\textwidth}{!}{\begin{threeparttable} \begin{tabular}{@{\extracolsep{5pt}} lrrr}  \toprule   & Chain & Grocery& Convenience\\  \midrule Mean store profits (conditional on remaining active) in 2010\$ & 73,074 & 42,719 & 43,937 \\  Mean entry costs (conditional on entering)  in 2010\$ & 129,169 & 192,093 & 244,170 \\  \\ \textit{Percentage change in mean store profits from} & & &  \\ \; One additional rival chain store (0-2 mi) & -9.69 & -30.45 & -16.01 \\  \; One additional rival chain store (2-5 mi) & -6.71 & -1.28 & -4.13 \\  \; One additional rival grocery store (0-2 mi) & -10.30 & -7.05 & -7.64 \\  \; One additional rival grocery store (2-5 mi) & -10.14 & -16.24 & -8.18 \\  \; One additional rival convenience store (0-2 mi) & -10.22 & -24.78 & -13.22 \\  \; One additional rival convenience store (2-5 mi) & 3.69 & -2.90 & -1.09 \\  \; One additional own chain store (0-2 mi) & -13.15 &  &  \\  \; One additional own chain store (2-5 mi) & 10.67 &  &  \\  \; Increase in dist. to distribution center by one s.d. from mean & -6.13 &  &  \\  \bottomrule  \end{tabular} \begin{tablenotes}
                \small
                \item \textit{Note: Averages are over all incumbent stores (for profits) and entrants (for entry costs) over the period 2010-2019. Conditional profits and entry costs include the expectation of the structural shock. Percentage changes are relative to the monopoly case. The mean distance to the closest distribution center is 190mi and the standard deviation is 130mi.} 
              \end{tablenotes}
              \end{threeparttable}} 